\documentclass[a4paper,10pt]{article}

\usepackage[T1]{fontenc}
\usepackage{hyperref}

%opening
\title{Blender JA Mixamo Importer Manual}
\author{BE-Arbiter}

\begin{document}

 \maketitle

 \newpage

 \tableofcontents

 \newpage

 \section{Purpose of this document}

 This manual explains the Blender add-on retargeting Mixamo animations (\texttt{.fbx}) onto Jedi
 Academy's \texttt{\_humanoid} skeleton as NLA strips. It is meant to be used together with the
 Blender Jedi Academy Plugin Suite (the \texttt{jediacademy} add-on), which reads and writes the
 game's \texttt{.gla} and \texttt{animation.cfg} files, and follows the same conventions as the
 JA dotXSI add-on. It does not explain how to use Blender itself.

 \section{Changelog}

 \subsection*{next version}
 Initial release:
 \begin{itemize}
  \item Mixamo FBX import (binary and ASCII) retargeted onto the \texttt{skeleton\_root} armature,
        as NLA strips or as the active action
  \item Automatic scale from both skeletons' hip heights, origin offset, in place and frame rate options
  \item Several sequence names separated by dots create stacked strips, one action per name
 \end{itemize}

 \section{Installation}

 Download \texttt{Blender-JA-Mixamo-Importer.zip} from the releases page and install it with the
 ``Install from Disk'' button in Blender's add-on preferences, then enable
 ``JA Mixamo NLA Import (.fbx)''. Blender 4.1 or newer is required.

 The add-on adds one menu entry: File $\rightarrow$ Import $\rightarrow$ JA Mixamo NLA Import (.fbx).

 \section{The Mixamo files}

 Download the animation from \url{https://www.mixamo.com} in the FBX format (``FBX Binary'' or
 ``FBX ASCII'', with or without skin, any frame rate). The file is read by the add-on itself:
 Blender's own FBX importer is not used (it refuses ASCII files).

 Only the skeleton and its animation are read. The bones are recognized by their Mixamo names,
 whatever their namespace (\texttt{mixamorig:Hips}, \texttt{mixamorig1:Hips}, \ldots); meshes and
 other models are ignored. The rest pose is the skeleton's own T-pose, as stored in the file.

 \section{Import}

 \subsection{Target armature}

 The animation is retargeted onto the scene's \texttt{skeleton\_root} armature (the name the
 jediacademy add-on gives it), or the active armature. Import \texttt{\_humanoid.gla}, or a
 \texttt{.glm} model using it, with the jediacademy add-on first. The bones must have the
 \texttt{\_humanoid} names.

 \subsection{Settings}

 \begin{description}
  \item[Scale] Converts the file's units for the pelvis motion. 0 (the default) measures it: the
        ratio between the armature's and the file's hip heights above the ankles.
  \item[Origin Offset] Translation added to the whole skeleton in every frame. Every stock
        \texttt{\_humanoid} animation has \texttt{model\_root} at (0, 0, -24): the player origin is
        24 units above the feet.
  \item[In Place] Drops the horizontal motion of the hips, keeping the vertical one.
  \item[FPS] Frame rate the animation is resampled to, 0 keeps the file's (30 for Mixamo).
  \item[Add as NLA Strip] Puts the animation in its own action on a new NLA strip, like the
        jediacademy add-on's \texttt{.gla} import with an \texttt{animation.cfg}. Otherwise the
        action becomes the armature's active action.
  \item[Strip Start] Frame of the new strip, -1 places it right after the last existing strip.
  \item[Sequence Name] Name of the action and of the \texttt{animation.cfg} sequence, the file name
        by default. Several names separated by dots (e.g. \texttt{BOTH\_A.BOTH\_B}) create one action
        per name, their strips stacked at the same frames on successive tracks.
  \item[Loop Frame] \texttt{animation.cfg} loop frame: -1 for no loop, 0 to loop from the start.
  \item[Set Scene Range \& FPS] Makes the scene's frame range include the new animation and uses
        its frame rate.
 \end{description}

 \subsection{Retargeting}

 Mixamo skeletons rest in a T-pose, \texttt{\_humanoid} in an A-pose. For every mapped bone the
 add-on first turns the Jedi Academy bone so that it points where the Mixamo bone points at rest,
 then applies the Mixamo bone's rotation in every frame. Only the pelvis takes the file's
 translation; every other bone keeps its length.

 \begin{tabular}{ll}
  Mixamo & \texttt{\_humanoid} \\
  \hline
  Hips & pelvis \\
  Spine, Spine1, Spine2 & lower\_lumbar, upper\_lumbar, thoracic \\
  Neck, Head & cervical, cranium \\
  Shoulder, Arm, ForeArm, Hand & clavical, humerus, radius, hand \\
  HandThumb1, 2 & d1\_j1, d1\_j2 \\
  HandIndex1, 2 & d2\_j1, d2\_j2 \\
  HandRing1, 2 & d4\_j1, d4\_j2 \\
  Hand & hang\_tag\_bone \\
  UpLeg, Leg, Foot & femurYZ, tibia, talus \\
 \end{tabular}

 \medskip

 The other bones (twist bones \texttt{humerusX}, \texttt{radiusX}, \texttt{femurX}, the face, the
 tails, \texttt{Motion}) follow their parent. The fingers and the hand tag bones follow the hand,
 although Ghoul 2 parents them to the forearm, so the saber stays in the hand.

 \section{Workflow with the Jedi Academy plugin suite}

 \begin{enumerate}
  \item Import \texttt{\_humanoid.gla} (or a model) with the jediacademy add-on, with its
        \texttt{animation.cfg} if the new animations are added to the existing ones.
  \item Import each Mixamo file with File $\rightarrow$ Import $\rightarrow$ JA Mixamo NLA Import
        (.fbx), giving it its sequence name (e.g. \texttt{BOTH\_STAND2}).
  \item Export the \texttt{.gla} and \texttt{animation.cfg} with the jediacademy add-on.
 \end{enumerate}

\end{document}
